\chapter{环形计数器——全部变式详解}

\section{变式1：N位环形计数器}

\begin{detailbox}[完整教学]
\textbf{【通用模板】}
\begin{lstlisting}
process(clk, rst_n)
begin
  if rst_n='0' then reg <= (N-1=>'1', others=>'0');  -- One-hot初值
  elsif rising_edge(clk) then
    reg <= reg(0) & reg(N-1 downto 1);  -- 右移+反馈
  end if;
end process;
\end{lstlisting}
\textbf{【结论】}N个DFF→N个状态。所有状态中恰好一个1（One-hot编码）。
\end{detailbox}


\section{变式2：约翰逊计数器（2N个状态）}

\begin{detailbox}[完整教学]
\textbf{【反馈】}\texttt{reg <= (not reg(0)) \& reg(N-1 downto 1);}
\textbf{【4位约翰逊状态序列】}$0000\to1000\to1100\to1110\to1111\to0111\to0011\to0001\to0000$
\textbf{【结论】}N个DFF→2N个状态（比环形计数器多一倍）。
\end{detailbox}


\section{变式3：自启动环形计数器}

\begin{detailbox}[完整教学]
\textbf{【问题】}干扰可能使环形计数器进入非法状态。
\textbf{【方案】}检测非法状态（如全0或多个1），强制恢复到合法状态(1000)。
\textbf{【VHDL】}在移位前检查：\texttt{if count\_ones(reg) /= 1 then reg <= "1000";}
\end{detailbox}


\section{变式4：用环形计数器实现时序控制器}

\begin{detailbox}[完整教学]
\textbf{【应用】}多周期操作器的步骤控制。
\textbf{【5步控制器】}5位环形计数器→每拍恰好一个步骤信号为1→控制对应步骤。
\textbf{【优点】}无需额外的译码逻辑——直接读Q值。
\end{detailbox}
